\documentclass[10pt]{article}
\usepackage[margin=1in]{geometry}
\usepackage{amsmath}
\usepackage{amssymb}
\usepackage{bbm}
\usepackage{booktabs}
\usepackage{float}
\usepackage{hyperref}
\usepackage{enumitem}

\title{Counterfactual Auction Computation}
\author{Konan Hara}
\date{October 4, 2026}

\begin{document}
\maketitle

\section{Purpose and Current Status}

This note documents the production counterfactual model and its Julia
implementation.  The production model replaces the earlier continuous-type
indirect-inference calculation with two discrete risk-aversion types and a
common upper bid support.  Interior bidding strategies solve a joint
inverse-strategy differential system parameterized by the lower-bid
probability.  Entry follows two
type-specific cost cutoffs.  The common-support restriction selects a single
equilibrium branch and avoids alternating between a flexible bid distribution
and an entry approximation.

The economic objects retained from the preceding implementation are the
effective-cost interpretation, the wholesale-only outside option, endogenous
portfolio choice on $[0,1]$, full-project awards, the objective-capacity
distribution, and policy comparisons with fixed-share auctions.  Capacity is
independent of $(\gamma,c)$ in the counterfactual type distribution.  The
factual distribution used to calibrate the registered pool continues to
preserve the estimated risk-aversion--cost correlation.

The final-precision production run completed on October 4, 2026.  It
solved all eight primitive specifications, eleven cost shifts, bidder share
choice, twenty-one fixed shares, and the matched-share comparison.  All 2,024
distinct equilibria passed the strict pooled-boundary and differential-system
checks.  The numerical values below are the canonical paper results.

\section{Counterfactual Primitives}

\subsection{Effective Cost and the Outside Option}

Let $\pi_{0i}$ denote bidder $i$'s outside certainty equivalent.  The auction
bids identify the effective cost
\begin{equation}
c_i=c_i^{physical}+\pi_{0i},
\label{eq:effective_cost}
\end{equation}
not its two components separately.  Estimation can therefore remain agnostic
about the specification of $\pi_{0i}$.  In the counterfactual, a project may
build entirely in the wholesale market, so
\begin{equation}
\pi_{0i}
=\max\left\{0,\mu_r-c_i^{physical}
-\frac{\gamma_i\sigma_r^2}{2}\right\}.
\end{equation}
A downward physical-cost shift $\Delta c\leq0$ maps the estimated effective
cost into
\begin{equation}
c_i(\Delta c)
=\max\left\{c_i(0)+\Delta c,
\mu_r-\frac{\gamma_i\sigma_r^2}{2}\right\}.
\label{eq:cost_floor}
\end{equation}
The lower bound is the effective cost of a wholesale-only project.  A project
at the bound is viable outside the auction and is counted in total built
capacity if it does not win auction support.

\subsection{Two Risk-Aversion Types}

The production approximation replaces the continuous distribution of
$\log\gamma$ with two equally weighted points.  If
$s_\gamma=\operatorname{SD}(\log\gamma)$, define
\begin{equation}
\log\gamma_L=\ell-d,
\qquad
\log\gamma_H=\ell+d,
\qquad
d=s_\gamma\sqrt{\frac{N-1}{N}},
\end{equation}
where $\ell$ is chosen so that the discrete distribution has the maintained
level mean $E[\gamma]=0.56427$.  The small finite-sample factor makes the
simulated sample standard deviation equal its target.  The production grid is
\begin{equation}
s_\gamma\in\{0.038337,0.088437\}.
\label{eq:gamma_grid}
\end{equation}
The first value is the point estimate.  The second corresponds to
$\operatorname{SD}(\gamma)=0.05$.  The more extreme bootstrap upper bound is
excluded from the production grid because it produced unstable winner
composition over adjacent fixed-share values under the common-support
approximation.

Conditional cost distributions are generated so the level correlation
between $\gamma$ and cost equals
\begin{equation}
\rho\in\{-0.658,0,0.132847,0.454\}.
\label{eq:rho_grid}
\end{equation}
The primary value is the Table~5 point estimate.  The endpoints are its
central 95 percent interval.  Capacity is drawn independently from its
empirical participant-weighted marginal distribution.  Thus capacity affects
competition and allocation but, conditional on $(\gamma,c)$, not own entry.

Within risk class $g\in\{L,H\}$, define cost above the wholesale floor as
\begin{equation}
x=c-\left(\mu_r-\frac{\gamma_g\sigma_r^2}{2}\right)\geq0.
\label{eq:x_type}
\end{equation}
The code represents each conditional distribution $F_g(x)$ by a weighted
Gaussian kernel distribution.  This smoothing is numerical quadrature for a
deterministic cutoff, not a behavioral entry shock.

\subsection{Synthetic Auction and Allocation}

The representative auction has the median number of registered projects and
uses participant-weighted auction primitives.  To construct the objective
capacity $\widetilde D$, the code simulates baseline sealed-bid auctions,
fills procurement capacity in units of $q_iCapacity_i$, prorates the marginal
project, and converts the resulting procurement draw into installed capacity.
It normalizes each draw by the simulated auction's expected aggregate
participant capacity and rescales it to the representative auction.  In a
counterfactual auction, every lower-bidding competitor contributes its full
$Capacity_j$ toward procurement:
\begin{equation}
W(P)=\Pr\left\{
\sum_{j=1}^{N-1}Capacity_j\mathbbm{1}\{b_j<b\}<\widetilde D
\right\},
\label{eq:winprob}
\end{equation}
where $P$ is the probability that a potential competitor enters and submits a
lower bid.  Because capacity is independent of strategic type, the code
computes $W(P)$ and $W_P(P)$ from the exact binomial mixture over the number of
lower bidders.  Conditional procurement probabilities are integrated using
common capacity and $\widetilde D$ draws.  This replaces repeated simulation
of a nonparametric winning-probability curve.

Awards are made to whole projects.  Realized awarded capacity can therefore
overshoot $\widetilde D$.  The same full-award rule and common allocation
random numbers are used in every policy scenario.

\section{Bidder Payoffs and Portfolio Choice}

Utility is measured relative to the outside payoff.  A bidder who wins with
bid price $b$, purchase-agreement share $q$, and effective cost $c$ receives
the incremental certainty equivalent
\begin{equation}
m(q,b;c,\gamma)
=q\widetilde\delta b+(1-q)\mu_r-c
-\frac{\gamma\sigma_r^2}{2}(1-q)^2.
\label{eq:surplus}
\end{equation}
If the bidder loses, incremental utility is $u_i(0)=0$.  The certainty
equivalent of the auction lottery is therefore
\begin{equation}
V
=-\frac{1}{\gamma_i}\log\left[
1-W(b)\{1-\exp(-\gamma_i m)\}\right],
\label{eq:auction_value}
\end{equation}
where utility is normalized relative to the outside payoff.  Qualification
costs $\kappa Capacity_i$.  Since winning surplus and qualification cost both
scale with own capacity, own capacity cancels from the entry comparison.

In the bidder-share-choice auction,
\begin{equation}
q(b,\gamma)
=\left[
1-\frac{\mu_r^{share}-\widetilde\delta b}
{\gamma\sigma_r^2}
\right]_{0}^{1}.
\label{eq:q_choice}
\end{equation}
For a fixed-share auction, $q=\lambda$.  The case $\lambda=0$ is implemented
as an explicit subsidy auction; bidders compete in bid price and can win even
though the purchase-agreement share is zero.

\section{Common-Support Equilibrium}

\subsection{Inverse First-Order Conditions}

Let
\begin{equation}
K(b)=-\frac{W(b)}{W'(b)}>0.
\label{eq:K}
\end{equation}
For an interior bid, the bid-price first-order condition implies
\begin{equation}
s_g(b,K)
=\frac{1}{\gamma_g}
\log\left[1+\gamma_gK(b)m_b(b,\gamma_g)\right],
\label{eq:surplus_inverse}
\end{equation}
where $s_g$ is equilibrium winning surplus and $m_b$ is the derivative of the
gross winning certainty equivalent with respect to bid.  In the share-choice
auction $m_b=q_g(b)\widetilde\delta$ by the envelope theorem.  In every
fixed-share auction, including $\lambda=0$, $m_b=\widetilde\delta$.

Equations \eqref{eq:surplus}--\eqref{eq:surplus_inverse} yield the cost cutoff
above the outside floor for each risk type:
\begin{equation}
x_g(b,K)
=G_g(b)-s_g(b,K)
-\left(\mu_r-\frac{\gamma_g\sigma_r^2}{2}\right),
\label{eq:inverse_cost}
\end{equation}
where $G_g(b)$ is gross winning CE before cost.  The probability of an entry
and bid below $b$ is
\begin{equation}
P(b,K)=\sum_{g\in\{L,H\}}\pi_gF_g(x_g(b,K)).
\label{eq:lower_bid_probability}
\end{equation}

\subsection{Probability-Parameterized Differential System}

The allocation map gives $W=W(P)$.  Bid price is monotone in cost conditional
on risk class, and $P$ is a bid CDF, so $P$ must increase with $b$.  Rather
than asking a fixed-step ODE in bid space to reproduce this restriction, the
production code uses $P$ itself as the independent variable.  Combining
$K=-W/(W_PP')$ with equation~\eqref{eq:lower_bid_probability} gives
\begin{align}
\frac{db}{dP}
&=-\frac{W_P(P)K}{W(P)}>0,
\label{eq:bid_probability_ode}\\
\frac{dK}{dP}
&=\frac{1-P_b(b,K)\,db/dP}{P_K(b,K)}.
\label{eq:K_probability_ode}
\end{align}
The system starts at the common upper bid, where $P$ equals the aggregate
entry probability, and integrates downward to $P=0.0005$.  Fourth-order
Runge--Kutta steps jointly update $(b,\log K)$.  At every stage, a Newton
projection imposes equation~\eqref{eq:lower_bid_probability}; the step size is
bounded both in probability and in its implied bid change.  This makes bid
monotonicity a property of the computational coordinate rather than an
ex-post diagnostic.  The code reports
\begin{equation}
\max_P\left|P(b(P),K(P))-P\right|
\label{eq:probability_ode_residual}
\end{equation}
as the differential-system residual.  At every accepted node, equation
\eqref{eq:inverse_cost} recovers both type-specific cost cutoffs.  The
resulting monotone inverse strategies map each entrant draw into its bid and
portfolio share.

The common upper bid is a computational approximation.  In an unrestricted
two-type model the two risk classes could have different upper supports, and
boundary conditions would require a more elaborate segment-by-segment
solution.  Requiring a common upper support yields a single pooled lower-bid
probability at the boundary and makes the probability-parameterized system
operational.  This is the main simplifying restriction used to retain
risk-based selection while
eliminating the unstable high-dimensional fixed point.

\subsection{Boundary Calibration and Continuation}

In each baseline sealed-bid auction, the code approximates the marginal
participant's bid by the simulated expectation of the maximum of $N_a$ draws
from that auction's fitted participant bid distribution, where $N_a$ is its
number of participants.  The participant-count-weighted mean of these
auction-specific values anchors the representative auction's common upper bid
at $\bar b=56.76584$.  For a candidate registered-project distribution, the
code solves $\bar K=-W(\bar b)/W'(\bar b)$ so that
equation~\eqref{eq:lower_bid_probability} at $(\bar b,\bar K)$ reproduces the
observed aggregate participation rate.  Let $V_g$ denote the CARA certainty
equivalent of the auction lottery for marginal type $g$.  The qualification
cost calibration is
\begin{equation}
\log\kappa=\sum_g\pi_g\log V_g.
\label{eq:kappa}
\end{equation}
For the primary primitive specification, this gives the representative
baseline calibration of $\kappa$.  Each correlation and dispersion sensitivity
specification is treated in turn as the baseline primitive model and receives
its own baseline calibration of $\kappa$ at $\Delta c=0$.
Within a primitive specification, the resulting $\kappa$ is held fixed across
all cost shifts and auction formats.  Thus policy comparisons never
recalibrate entry cost, while sensitivity comparisons allow the nuisance
entry-cost calibration to remain consistent with the observed factual entry
rate under the alternative primitives.

In a counterfactual scenario, the two boundary variables $(\bar b,\bar K)$
are required to satisfy the pooled marginal-value condition
\begin{equation}
\sum_g\pi_g\log V_g(\bar b,\bar K)=\log\kappa.
\label{eq:pooled_entry}
\end{equation}
One equation does not identify both boundary variables.  The implementation
selects the point on this level set closest to the translated preceding
solution:
\begin{equation}
\left(\frac{\bar b-\bar b^{pred}}{2}\right)^2
+\left(\frac{\log\bar K-\log\bar K^{pred}}{0.5}\right)^2.
\label{eq:continuation_metric}
\end{equation}
Cost shifts are solved sequentially at
$\Delta c=0,-1,-2,\ldots,-10$.  The preceding cost solution supplies the next
prediction.  This continuation metric is a numerical branch-selection rule,
not an additional behavioral restriction.

Under common support, equation~\eqref{eq:pooled_entry} need not make both
type-specific marginal values equal $\kappa$.  The code therefore reports
\begin{equation}
r_g=\log(V_g/\kappa),\qquad
R_{RMSE}=\left(\frac{r_L^2+r_H^2}{2}\right)^{1/2},\qquad
R_{max}=\max_g|r_g|.
\label{eq:type_residual}
\end{equation}
These are approximation diagnostics.  The pooled residual in
equation~\eqref{eq:pooled_entry}, the residual in
equation~\eqref{eq:probability_ode_residual}, monotonicity of $b(P)$, and
positivity of $K$ are strict numerical checks.

\section{Recovering the Registered Pool}

The primitive participant distribution is conditional on participation,
while participation must be evaluated over registered projects.  For project
draw $j$, define
\begin{equation}
w_j(\alpha)
=\frac{\exp\{\alpha'h(\gamma_j,c_j)\}}
{\sum_k\exp\{\alpha'h(\gamma_k,c_k)\}},
\label{eq:registered_weight}
\end{equation}
where $h(\gamma,c)$ contains standardized $\log\gamma$ and cost, their centered
squares, and their interaction.  These weights describe a candidate
registered-project distribution.  For each candidate $\alpha$, the code holds
the baseline marginal-participant bid anchor $\bar b$ fixed, solves $\bar K$
so that the type-specific cutoffs $\bar x_g(\bar b,\bar K)$ reproduce the
observed participant--registrant ratio, and applies those cutoffs to the
weighted draws.  It calibrates $\alpha$ so that the selected draws reproduce
the primitive participant moments.  Thus the procedure inverts the model's
participation selection: it recovers the registered-project weights that
generate the primitive participant distribution after the baseline
participation cutoffs are applied.

The common boundary then implies the type-specific marginal auction values
used in equation~\eqref{eq:kappa}.  Each correlation and dispersion sensitivity
is an alternative baseline calibration and receives its own registered-pool
weights, $\bar K$, and $\kappa$.  Within a primitive scenario, the recovered
registered weights and calibrated $\kappa$ are fixed over all cost reductions
and auction mechanisms.

\section{Policy Outcomes}

The production calculation compares:
\begin{enumerate}[leftmargin=2em]
\item bidder share choice;
\item fixed shares $\lambda=0,0.05,\ldots,1$; and
\item a fixed share equal to the capacity-weighted winner mean share in the
corresponding share-choice auction.
\end{enumerate}
For every mechanism and cost shift, the code reports entry, entrant and winner
mean cost, mean risk aversion, the high- and low-risk-type shares, bids, markups, winner
mean $q$, auction-supported capacity, capacity built without auction support, and total built
capacity.  The policymaker outcomes are expected net cost and
\begin{equation}
CECost_{PM}=E[Cost_{PM}]
+\frac{\gamma_{PM}}{2}\operatorname{Var}(Cost_{PM}),
\end{equation}
for $\gamma_{PM}=0$ and $\gamma_{PM}=E[\gamma]$.
The simulation constructs this object per MWh of objective generation.  The
paper figures report aggregate policymaker CE in billion dollars by multiplying
it by $\widetilde D\times 8{,}760\times20$.  Equivalently, this multiplies a
winner-generation CE net of the small full-award capacity overshoot by total
winner generation.  The matched-share table reports policymaker CE and the
capacity-weighted winner markup using the same conversion, in \$100 million.

The matched-share comparison holds average realized public risk exposure
approximately fixed, but it is not a dominance test.  The matched fixed-share
auction uses the average share that bidders reveal under endogenous portfolio
choice and therefore does not preserve their informational advantage from
choosing $q$ individually.  Moreover, total built capacity can differ across
the two mechanisms.  Results should be described as combinations of policy
cost, risk exposure, and construction rather than as a one-dimensional welfare
ranking.

\section{Production Grid and Computation}

The final-precision settings are
\begin{center}
\begin{tabular}{lr}
\toprule
Object & Draws\\
\midrule
Auxiliary primitive draws & 5,000\\
Winning-probability/allocation-map simulations & 1,000\\
Allocation simulations & 2,000\\
$\widetilde D$ draws & 2,000\\
\bottomrule
\end{tabular}
\end{center}
The Cartesian product of equations~\eqref{eq:gamma_grid} and
\eqref{eq:rho_grid} gives eight independent primitive scenarios.  Each task
solves the complete cost and policy grid so cost continuation remains internal
to a process.  The local production command is
\begin{verbatim}
CF_FRESH=1 CF_WORKERS=4 ./counterfactual.sh
\end{verbatim}
This runs eight scenarios locally with at most four Julia processes.  Each
Julia process itself uses one thread.

\section{Files and Audit Checks}

The live implementation contains the following files:
\begin{itemize}[leftmargin=2em]
\item \path{counterfactual_model.jl}: two-type primitives, exact allocation
map, inverse FOCs, boundary solver, probability-parameterized system, and
registered-pool calibration;
\item \path{counterfactual_prep.jl}: representative-auction primitives,
objective-capacity simulation, and the baseline marginal-participant bid
anchor;
\item \path{counterfactual_baseline.jl}: baseline registered-pool, $\overline
K$, and $\kappa$ calibration;
\item \path{counterfactual.jl}: one complete primitive-scenario task;
\item \path{counterfactual_result.jl}: aggregation and policy comparisons;
\item \path{counterfactual.sh}: local production entry point; and
\item the three short shell wrappers used by that entry point.
\end{itemize}

\begin{sloppypar}
Large task files remain under
\path{Dropbox/renewable-auction/intermediate/counterfactual_task}.  Aggregated
CSV files are written to
\path{Dropbox/renewable-auction/intermediate/counterfactual_output}.  Logs are
written to \path{Dropbox/renewable-auction/log/counterfactual}.  The script
\path{code/result/counterfactual.jl} creates compact repository outputs.  The
production result wrapper then runs \path{code/result/counterfactual.R} to
create the paper tables, figures, and sentence files.
\end{sloppypar}

\section{Final Production Results}

The primary share-choice results are:
\begin{center}
\begin{tabular}{lrrrrrr}
\toprule
$\Delta c$ & 0 & $-2$ & $-4$ & $-6$ & $-8$ & $-10$\\
\midrule
Entry & 0.668 & 0.668 & 0.668 & 0.668 & 0.668 & 0.668\\
Winner mean $q$ & 1.000 & 0.993 & 0.868 & 0.677 & 0.547 & 0.482\\
High-$\gamma$ winner share & 0.399 & 0.399 & 0.393 & 0.427 & 0.807 & 0.986\\
Capacity without support (MW) & 0.0 & 0.0 & 0.0 & 19.4 & 532.3 & 2,082.5\\
Total built capacity (MW) & 486.3 & 486.4 & 486.5 & 505.9 & 1,018.9 & 2,569.1\\
\bottomrule
\end{tabular}
\end{center}
The capacity-weighted winner share decreases monotonically in every primitive
specification, and capacity built without auction support increases monotonically.  Entry is
nearly constant.  Across the primary policy and cost grid it ranges from
0.668 to 0.669, and the largest adjacent share-choice change across the full
sensitivity grid is 0.019 percentage point.  Equilibrium bid adjustments keep
the marginal auction value near the fixed qualification threshold, stabilizing
the number of competitors.  Stable entry does not mechanically fix markups
because the bid first-order condition changes with risk sharing; empirically,
however, the capacity-weighted winner markup across bidder share choice and
fixed shares $\lambda\in\{0,0.5,1\}$ ranges only from \$0.16/MWh to
\$1.32/MWh.  The economic phase-out therefore operates through
portfolio choice and wholesale construction rather than a collapse in auction
participation.

Selection responds both to policy and to the outside option.  At
$\Delta c=-8$, high-$\gamma$ projects account for 71.9 percent of subsidy-auction
winners and 85.9 percent of full-purchase-agreement winners, so leaving more
wholesale risk with investors selects more strongly toward low-risk-premium
projects.  Along the cost path, however, low-$\gamma$ projects become
wholesale-viable first and increasingly build without auction support.  This correlated cost
selection drives the high-$\gamma$ winner share up sharply in the deepest
cost reductions even though the entrant pool remains approximately balanced.

For a risk-neutral policymaker, the best fixed share is $\lambda=1$ at every
cost shift.  For a policymaker with $\gamma_{PM}=E[\gamma]$, the best fixed
share is $\lambda=0.5$.  In the matched-share comparison, the fixed-share
auction generally has slightly lower policymaker certainty-equivalent cost
and slightly more total construction.  This is not a dominance result because
the fixed policy uses the average share revealed under endogenous portfolio
choice and therefore does not preserve bidders' informational advantage.

\section{Final Numerical Audit}

The completed aggregation contains 4,048 rows: 2,024 distinct equilibria and
two policymaker-risk evaluations for each.  There are no duplicate scenario
keys and no missing core outcomes.  All 2,024 equilibria satisfy bid
monotonicity, probability monotonicity, and $K>0$, and every allocation meets
objective capacity.  The maximum pooled marginal-value residual is
$1.15\times10^{-7}$ and the maximum probability-system residual is below
$10^{-10}$.

The common-support restriction is the material approximation.  In the
primary specification, $R_{max}$ is at most 0.044.  It reaches 0.117 in the
larger-dispersion, negatively correlated specification.  Fixed-share outcome
paths remain smooth: across the full grid, the largest adjacent entry change
over a 0.05 change in $\lambda$ is 0.071 percentage point, the largest adjacent
winner-composition change is 6.4 percentage points, and the largest reversal
is 0.24 percentage point.  Some boundary $K$ values move sharply despite these
smooth observables, so economic outcomes are better conditioned than the
individual boundary parameters.

The primary baseline calibration gives $\bar b=56.76584$ and
$\overline K=8.17775$, reproduces the observed participation rate 0.667692,
and matches registered-pool moments with RMSE $8.94\times10^{-5}$.  The
calibrated $\kappa$ is approximately $4.99\times10^{-19}$ per MWh.  This small
value reflects the extremely low
winning probability of the marginal entrant under the common-upper-support
approximation and therefore depends on how the winning-probability function is
extrapolated in the far tail.  Since counterfactual bidding depends on the
winning-probability schedule and its slope through $K=-W/W'$, the magnitude of
$\kappa$ should not be interpreted independently.  As long as the tail winning
probability remains positive and large enough to affect the marginal bidder's
bidding behavior, changing its precise extrapolation and the implied scale of
$\kappa$ should not materially alter the economic results.

\end{document}
